\documentclass[10pt,a4paper]{amsart}
\usepackage[margin=19mm]{geometry}
\usepackage{amsmath,amssymb,mathtools}
\usepackage[scr=rsfs,cal=euler]{mathalfa}
\usepackage{xfrac}
\usepackage[unicode,hidelinks,bookmarks=false]{hyperref}
\newcommand{\ie}{i.e.\ }
\newcommand{\cf}{cf.\ }
\newcommand{\resp}{respectively\ }
\newcommand{\Subsection}{\subsection}
\providecommand{\nobreakdash}{\nobreak\hbox{-}}
\setlength{\emergencystretch}{2em}
\multlinegap0pt
\usepackage{bm}
\usepackage{bbm}
\usepackage{dsfont}
\usepackage{xspace}
\usepackage{cancel}
\usepackage{mathtools}
\usepackage{amsthm}
\makeatletter
\@ifundefined{thm}{\newtheorem{thm}{Theorem}}{}
\makeatother
\title[Trees with non log-concave independent set sequences]{Trees with non log-concave independent set sequences}
\author{David Galvin}
\date{Source: arXiv:2502.10654v2 (CC BY 4.0)}
\begin{document}
\maketitle

\noindent\emph{Source and licence.} Trees with non log-concave independent set sequences. Version arXiv:2502.10654v2 (CC BY 4.0), distributed under the Creative Commons Attribution 4.0 International licence, as recorded in the arXiv licence metadata for this exact version and shown on the abstract page https://arxiv.org/abs/2502.10654v2. Full rights notice: licenses/arxiv-2502.10654-CC-BY-4.0.txt.\par\smallskip

\noindent\textit{Editorial scope.} Counted unit: the Thm whose source label is \texttt{thm-more-general} in the article (source0.tex lines 129--134, source label \texttt{thm-more-general}), delivered together with the article's own complete original proof of it (source0.tex lines 140--202, one \texttt{proof} environment of 63 lines). No mathematics is written, repaired or re-derived: statement and proof are the article's own text line for line. The proof is detached from the statement in the source: the article states the result at lines 129--134 and prints its proof at lines 140--202, and the proof environment's own header names the statement, so the association is the article's own. Required context retained uncounted: none. Every definition, display and label of those lines is retained unchanged. Omitted from this package and not redistributed: 1--128, 135--139, 203--350. Nothing inside a retained excerpt is omitted or altered: every retained body is delivered exactly as the article's lines are written, so no omission receipt of any kind accompanies this package. Citations: the retained excerpts carry no citation command at all, so no bibliography entry of the article is transcribed and no reference is redistributed. Reference adaptation: none was needed and none was made; every reference inside the delivered unit resolves inside it, so no unresolved reference or citation is produced. Printed slips kept literal: no printed word, symbol, hypothesis or number of the counted statement or of its proof was altered. Layout and class: the article's own class is \texttt{article} with packages including latexsym, amsmath, amsthm, amssymb, graphicx, appendix, tikz, hyperref, setspace, geometry, url, xcolor; the wrapper is the lane's pinned \texttt{amsart} editorial wrapper at 10pt on A4, which restates the theorem-like environments the retained text needs and adds the article's own macro definitions that the retained text needs (none), with extra wrapper packages (bm, bbm, dsfont, xspace, cancel, mathtools). The delivered numbering is the unit's own. Assets: no retained span contains a figure environment, a graphics-inclusion command, a PDF-page inclusion, a table or a TikZ picture; no image file is copied into this package, the package declares no asset dependency and no image file is redistributed with it. Licence: version arXiv:2502.10654v2 of this article is distributed under the Creative Commons Attribution 4.0 International licence, as recorded in the arXiv licence metadata for this exact version and shown on the abstract page https://arxiv.org/abs/2502.10654v2. A full rights notice is delivered at licenses/arxiv-2502.10654-CC-BY-4.0.txt. Characters: every retained excerpt is pure ASCII, so the delivered engine, XeLaTeX with its default Latin Modern text font, reports no missing character.
\begin{thm} \label{thm-more-general}
For $t \leq m = m(t) \leq 2^{t/16}$ we have that for large enough $t$ the log-concavity of the independent set sequence of $T_{m,t,1}$ is broken at $mt+2$; that is,
$$
\left(i_{T_{m,t,1}}(mt+2)\right)^2 < i_{T_{m,t,1}}(mt+1)i_{T_{m,t,1}}(mt+3). 
$$
\end{thm}

\begin{proof} (Proof of Theorem \ref{thm-more-general})
Denote by $f_{m,t}(k)$ the number of independent sets in $T_{m,t,1}$ that do not include the root $v$, and by $h_{m,t}(k)$ the number that do include $v$. Observe that on deleting $v$ and all its neighbors $T_{m,t,1}$ reduces to a perfect matching on $mt$ edges, a graph that admits independent sets of size only up to $mt$. It follows that $h_{m,t}(k)=0$ for $k \geq mt+2$, and that 
\begin{equation} \label{split}
i_{T_{m,t,1}}(k) = \left\{
\begin{array}{rl}
h_{m,t}(k) + f_{m,t}(k) & \mbox{if $k \leq mt+1$} \\
f_{m,t}(k) & \mbox{if $k \geq mt+2$}.
\end{array}
\right.
\end{equation}
Further, we have $h_{m,t}(mt+1)=2^{mt}$. So our goal is to show that
$$
\left(f_{m,t}(mt+2)\right)^2 < (2^{mt} + f_{m,t}(mt+1))f_{m,t}(mt+3),
$$
which is implied by 
\begin{equation} \label{mt}
\left(f_{m,t}(mt+2)\right)^2 < 2^{mt}f_{m,t}(mt+3).
\end{equation}
We will establish \eqref{mt} by showing that for $t \leq m \leq 2^{t/16}$ we have
\begin{equation} \label{inq-f2,f3}
f_{m,t}(mt+3) \geq \binom{m}{3}2^{mt-3t}~~~\mbox{and}~~~f_{m,t}(mt+2) = \left(1+o(1)\right) \binom{m}{2}2^{mt-2t}.
\end{equation}
That \eqref{inq-f2,f3} implies \eqref{mt} for large $t$ and $m$  satisfying $m \leq 2^{t-1}$ (not just $m \leq 2^{t/16}$) is easily established.  

\smallskip

We get the lower bound on $f_{m,t}(t^2+3)$ in \eqref{inq-f2,f3} by considering those independent sets of size $mt+3$ that include exactly $3$ vertices from $\{w_1, \ldots, w_m\}$ (and so necessarily don't contain the root). There are $\binom{m}{3}$ ways to select the three vertices. Once they are chosen the remaining $mt$ vertices come from a graph that consists of $3t$ isolated vertices and $mt-3t$ isolated edges, for which there are $2^{mt-3t}$ options.

For the asymptotic determination of $f_{m,t}(mt+2)$ in \eqref{inq-f2,f3}: 
\begin{itemize}
\item Arguing exactly as we did for $f_{m,t}(mt+3)$, there are $\binom{m}{2}2^{mt-2t}$ independent sets in $T_{m,t,1}$ of size $mt+2$ that include exactly $2$ vertices from $\{w_1, \ldots, w_m\}$ (and so necessarily don't contain the root). 
\item Next, there is no contribution to $f_{m,t}(mt+2)$ from independent sets that include exactly $0$ or $1$ vertices from $\{w_1, \ldots, w_m\}$ (the maximum size of such independent sets is $mt+1$).
\item The remaining independent sets that contribute to $f_{m,t}(t^2+2)$ have $s$ vertices from $\{w_1, \ldots, w_m\}$, where $s$ satisfies $3 \leq s \leq m$ ($\binom{m}{s}$ options). The remaining $mt-(s-2)$ vertices come from a graph that consists of $st$ isolated vertices and $mt-st$ independent edges. To construct all such independent sets, we can first choose $\ell$ of the isolated edges from which not to select a vertex, then choose a vertex from each of the remaining edges, and finally choose $(s-2)-\ell$ isolated vertices to not select to complete the independent set. Here $\ell$ satisfies $0 \leq \ell \leq \min\{s-2,mt-st\}$.   
\end{itemize}
It follows from the discussion above that for each $3 \leq s \leq m$ and $0 \leq \ell \leq \min\{s-2,mt-st\}$ we have a term
$$
\binom{m}{s}\binom{mt-st}{\ell}\binom{st}{(s-2)-\ell}2^{mt-st-\ell}
$$
contributing to $f_{m,t}(mt+2)$, and that these are the only terms contributing to $f_{m,t}(mt+2)$ aside from $\binom{m}{2}2^{mt-2t}$ (which we hope to show is the dominant term). Since there at most $m^2$ such terms, we complete the verification of the asymptotic expression for $f_{m,t}(mt+2)$ in \eqref{inq-f2,f3} by showing that
$$
\frac{\binom{m}{s}\binom{mt-st}{\ell}\binom{st}{(s-2)-\ell}2^{mt-st-\ell}}{\binom{m}{2}2^{mt-2t}} = o\left(m^{-2})\right)
$$
or equiavlently
\begin{equation} \label{asy}
\frac{\binom{m}{s}\binom{mt-st}{\ell}\binom{st}{(s-2)-\ell}}{2^{(s-2)t+\ell}} = o(1).
\end{equation}
Using (here and throughout) the simple bound $\binom{n}{k} \leq n^k$ we have
\begin{eqnarray*}
\binom{m}{s} & \leq & m^s = 2^{s \log m}, \\
\binom{mt-st}{\ell} & \leq & (mt)^s \leq 2^{2s\log m}
\end{eqnarray*}
(the latter using $\ell \leq s$ and $t \leq m$) and (for $t \geq 2$, and using $s,t \leq m$)
$$
\binom{st}{(s-2)-\ell} \leq (st)^s \leq 2^{2s\log m}.
$$
Lower bounding $(s-2)t+\ell \geq st/3$ (recall $s \geq 3$) we get that the left-hand side of \eqref{asy} is upper bounded by 
$$
2^{5s\log m-(st)/3}
$$ 
As long as $m \leq 2^{t/16}$ this is $o(1)$.
 
This concludes the verification that $f_{m,t}(mt+2)=(1+o(1))\binom{m}{2}2^{mt-2t}$, and so the verification that the log-concavity of $T_{m,t,1}$ is broken at $mt+2$ for sufficiently large $t$ as long as $m \leq 2^{t/16}$.
\end{proof}

\end{document}
